\documentclass[11pt]{article}

\usepackage[margin=0.75in]{geometry}
\usepackage{booktabs}
\usepackage{amsmath}
\usepackage{multirow}
\usepackage{hyperref}

\title{Source-Domain Retention: How Much Do Adapted Checkpoints Forget?}
\author{}
\date{\today}

\begin{document}
\maketitle

\section{Motivation and setup}

Every comparison elsewhere in this study scores a checkpoint on its \emph{target}
domain -- the domain the MMD term (or plain fine-tune) was meant to adapt \emph{to}.
This document asks the complementary question: after that adaptation/fine-tune,
how much of the \emph{source} domain does the checkpoint still recognize?

Every checkpoint below that reached a paper table via regime no\_mmd, naive\_mmd,
reg\_mmd, or reg\_mmd\_smaller\_weight (cross-referenced against
\texttt{experiment\_summary\_v6.tex}, \texttt{experiment\_summary\_v7.tex}, and
\texttt{freeze\_epoch\_ablation.tex}) was re-evaluated with a plain \texttt{model.val()}
call on its own source domain's held-out val split -- the same split used to select
\texttt{best.pt} during that domain's original pretraining run. Baseline/Solo
checkpoints are excluded throughout: they train on the target domain from a
COCO-direct start and never see a source domain, so retention does not apply to them.
Imgsz 1920 throughout, matching training resolution.

\paragraph{A caveat.} These tables report only the \emph{post-adaptation} source-domain
score. The natural reference point -- how well the original pretrained-on-source
checkpoint scored on this same val split, before it was ever repurposed for
adaptation -- is not yet included here; see Section~\ref{sec:open}.

\section{Real$\to$Syn direction: source = real\_large}

Source = real\_large, target = syn\_small. \texttt{reg\_mmd} and
\texttt{reg\_mmd\_smaller\_weight} are the \texttt{-4} reruns (matching
\texttt{experiment\_summary\_v7.tex}'s real$\to$syn table, which superseded the
earlier unsuffixed runs used in \texttt{experiment\_summary\_v6.tex}).

\begin{table}[h!]
    \centering
    \begin{tabular}{lrrrr}
        \toprule
        Model & person mAP50 & person mAP50-95 & vehicle mAP50 & vehicle mAP50-95 \\
        \midrule
        no\_mmd                  & 0.020 & 0.007 & 0.350 & 0.213 \\
        naive\_mmd               & 0.025 & 0.006 & 0.307 & 0.191 \\
        reg\_mmd                 & 0.040 & 0.010 & 0.333 & 0.210 \\
        reg\_mmd\_smaller\_weight & 0.030 & 0.007 & 0.319 & 0.199 \\
        \bottomrule
    \end{tabular}
    \caption{Evaluated on real\_large val (548 images) -- i.e.\ how well each
    syn\_small-adapted checkpoint still detects the real domain it started from.}
\end{table}

\noindent All four variants collapse far below the $\sim$0.67--0.70 person / $\sim$0.69--0.72
vehicle mAP50 this same source domain reaches when evaluated on itself right after
pretraining (Table~1 of \texttt{experiment\_summary\_v6.tex}, target-domain column).
Both MMD variants retain slightly more than no\_mmd on every column, and reg\_mmd is
the best of the three MMD variants here -- a small but consistent anchoring effect
from the source-domain batches seen during adaptation.

\section{ClearNoon source: syn\_large\_clearnoon}

Source = syn\_large\_clearnoon, target = real\_small. Bandwidth-freeze-epoch = 1 for
the regularized variants (matching the ClearNoon side of
\texttt{experiment\_summary\_v6.tex} Section~2.1).

\begin{table}[h!]
    \centering
    \begin{tabular}{lrrrr}
        \toprule
        Model & person mAP50 & person mAP50-95 & vehicle mAP50 & vehicle mAP50-95 \\
        \midrule
        no\_mmd                  & 0.085 & 0.020 & 0.273 & 0.173 \\
        naive\_mmd               & 0.108 & 0.026 & 0.287 & 0.177 \\
        reg\_mmd                 & 0.084 & 0.019 & 0.283 & 0.176 \\
        reg\_mmd\_smaller\_weight & 0.091 & 0.020 & 0.286 & 0.180 \\
        \bottomrule
    \end{tabular}
    \caption{Evaluated on syn\_large\_clearnoon val.}
\end{table}

\noindent Here naive\_mmd -- the unregularized variant -- retains the most source
signal on both classes, ahead of both regularized variants and no\_mmd. This is the
one group in this study where regularizing the MMD term does not help retention.

\section{All-weather synthetic source: syn\_large}

Source = syn\_large (6400 images, all four weather conditions), target = real\_small.
Both freeze-epoch settings are included since both appear in
\texttt{freeze\_epoch\_ablation.tex}.

\begin{table}[h!]
    \centering
    \small
    \begin{tabular}{lrrrr}
        \toprule
        Model & person mAP50 & person mAP50-95 & vehicle mAP50 & vehicle mAP50-95 \\
        \midrule
        no\_mmd                                    & 0.100 & 0.025 & 0.201 & 0.129 \\
        naive\_mmd                                 & 0.101 & 0.025 & 0.200 & 0.129 \\
        reg\_mmd, freeze=1                          & 0.112 & 0.026 & 0.230 & 0.148 \\
        reg\_mmd, freeze=5                          & 0.125 & 0.032 & 0.219 & 0.142 \\
        reg\_mmd\_smaller\_weight, freeze=1          & 0.120 & 0.030 & 0.227 & 0.146 \\
        reg\_mmd\_smaller\_weight, freeze=5          & 0.109 & 0.028 & 0.222 & 0.143 \\
        \bottomrule
    \end{tabular}
    \caption{Evaluated on syn\_large val (640 images).}
\end{table}

\noindent Both regularized variants clearly beat no\_mmd/naive\_mmd on every column
here -- unlike the ClearNoon group above, regularization helps retention when the
source domain is the full four-weather set. reg\_mmd at freeze=5 gives the best
person retention of the whole study (0.125/0.032); reg\_mmd at freeze=1 gives the
best vehicle retention (0.230/0.148). There is no single freeze-epoch setting that
wins on both classes at once, echoing the person/vehicle trade-off already documented
on the target-domain side in \texttt{experiment\_summary\_v6.tex} Section~2.2.

\section{Vehicle-only source: syn\_large\_vehicle\_only}

Source = syn\_large\_vehicle\_only, target = real\_small\_vehicle\_only (person
removed from training entirely). Bandwidth-freeze-epoch = 5.

\begin{table}[h!]
    \centering
    \begin{tabular}{lrr}
        \toprule
        Model & vehicle mAP50 & vehicle mAP50-95 \\
        \midrule
        no\_mmd                  & 0.211 & 0.137 \\
        naive\_mmd               & 0.216 & 0.140 \\
        reg\_mmd                 & 0.225 & 0.147 \\
        reg\_mmd\_smaller\_weight & 0.223 & 0.147 \\
        \bottomrule
    \end{tabular}
    \caption{Evaluated on syn\_large\_vehicle\_only val. No person column: person was
    never part of this source or target dataset.}
\end{table}

\noindent Both regularized variants beat no\_mmd and naive\_mmd here, with reg\_mmd
slightly ahead of reg\_mmd\_smaller\_weight -- the same direction as the all-weather
syn\_large group above, on the one class both groups share.

\section{syn\_xlarge as source}

Source = syn\_xlarge (10240 images, all weathers), across four different targets/
scales: real\_small (YOLOv10n and YOLOv10m) and real\_large at two bandwidth-freeze
settings.

\begin{table}[h!]
    \centering
    \small
    \begin{tabular}{lrrrr}
        \toprule
        Model & person mAP50 & person mAP50-95 & vehicle mAP50 & vehicle mAP50-95 \\
        \midrule
        $\to$real\_small, reg\_mmd (freeze=5)                     & 0.139 & 0.035 & 0.236 & 0.155 \\
        $\to$real\_small, reg\_mmd, YOLOv10m (freeze=5)           & 0.133 & 0.035 & 0.222 & 0.146 \\
        $\to$real\_large, reg\_mmd (freeze=6)                     & 0.104 & 0.027 & 0.219 & 0.142 \\
        $\to$real\_large, reg\_mmd, stronger weight (freeze=25)   & 0.138 & 0.036 & 0.238 & 0.157 \\
        \bottomrule
    \end{tabular}
    \caption{Evaluated on syn\_xlarge val. No no\_mmd/naive\_mmd control exists yet
    at this source scale (see Open items).}
\end{table}

\noindent Within the real\_large-target pair, the longer, later-freezing
stronger-weight variant (freeze=25, 120 epochs) retains source performance
noticeably better than the freeze=6 variant on every column -- the opposite of what
its already-narrower target-domain gap (\texttt{experiment\_summary\_v7.tex}
Section~3) might suggest: this variant is not just closer to solo on the target
side, it also forgets its source less. The YOLOv10n vs.\ YOLOv10m comparison
(both real\_small-target, freeze=5) shows the smaller backbone retaining slightly
more of the source domain on every column, despite YOLOv10m winning on the
target-domain side (\texttt{experiment\_summary\_v7.tex} Section~3.1) -- larger
capacity here appears to specialize toward the target at a small additional cost
to the source.

\section{Overall pattern}

Across all five source domains, one pattern holds without exception: every adapted
or no-MMD checkpoint's source-domain score is a small fraction of what that same
domain reaches when evaluated on itself immediately after pretraining -- adaptation
training, MMD-regularized or not, does not preserve source-domain competence.
Within that collapse, MMD generally retains more than plain no-MMD fine-tuning
(4 of 5 groups), and regularizing the MMD term generally retains more than the
unregularized naive variant (3 of 4 groups where both exist) -- but the ClearNoon
source group inverts both of these directions, so neither is a universal rule.

\section{Open items}
\label{sec:open}

\begin{itemize}
    \item No pre-adaptation reference point: each source domain's own pretrain
    checkpoint, evaluated on this same val split immediately after pretraining,
    would turn these absolute numbers into a proper retention \emph{delta}.
    \item No no\_mmd/naive\_mmd control at the syn\_xlarge source scale, for either
    the real\_small or real\_large target -- only regularized variants have been
    run so far at this source scale.
    \item person mAP50-95 in particular is often an order of magnitude below the
    vehicle column across every group; whether this reflects genuine source-domain
    person forgetting or a checkpoint-selection artifact (person's volatility
    already flagged in \texttt{experiment\_summary\_v6.tex} Section~2.2) has not
    been separately investigated here.
\end{itemize}

\end{document}
